\documentclass[11pt]{article}
\usepackage[utf8]{inputenc}
\usepackage{amsmath}
\usepackage{geometry}
\usepackage{listings}
\usepackage{xcolor}

\geometry{letterpaper, margin=1in}

\definecolor{codegreen}{rgb}{0,0.6,0}
\definecolor{codegray}{rgb}{0.5,0.5,0.5}
\definecolor{codepurple}{rgb}{0.58,0,0.82}
\definecolor{backcolour}{rgb}{0.95,0.95,0.92}

\lstdefinestyle{mystyle}{
    backgroundcolor=\color{backcolour},   
    commentstyle=\color{codegreen},
    keywordstyle=\color{magenta},
    numberstyle=\tiny\color{codegray},
    stringstyle=\color{codepurple},
    basicstyle=\ttfamily\footnotesize,
    breakatwhitespace=false,         
    breaklines=true,                 
    captionpos=b,                    
    keepspaces=true,                 
    numbers=left,                    
    numbersep=5pt,                  
    showspaces=false,                
    showstringspaces=false,
    showtabs=false,                  
    tabsize=2
}

\lstset{style=mystyle}

\title{\textbf{Editorial: Problem B - The 727 Jumpscare}}
\author{Setter: Hudson}
\date{}

\begin{document}

\maketitle

\section*{Problem Overview}
We are given a string $S$ consisting of digits. We need to find the number of subsequences "727". That is, the number of index triples $(i, j, k)$ such that $1 \le i < j < k \le |S|$ and $S[i] = '7', S[j] = '2', S[k] = '7'$.

\section*{Naive Solution}
A brute-force approach would be to iterate through all possible triples $(i, j, k)$ and check the condition.
This would require three nested loops:
\begin{lstlisting}[language=C++]
long long count = 0;
int n = s.length();
for (int i = 0; i < n; i++) {
    for (int j = i + 1; j < n; j++) {
        for (int k = j + 1; k < n; k++) {
            if (s[i] == '7' && s[j] == '2' && s[k] == '7') {
                count++;
            }
        }
    }
}
\end{lstlisting}
The time complexity is $O(N^3)$. With $|S| \le 5000$, $N^3 \approx 1.25 \times 10^{11}$ operations, which exceeds the operation limit for 1.0 second ($\approx 10^8$ operations). This solution will result in Time Limit Exceeded (TLE).

\section*{Optimal Solution (Dynamic Programming)}
We can solve this in a single pass $O(N)$ by keeping track of partial subsequences. The problem asks us to build "727". This is built in steps: "7" $\rightarrow$ "72" $\rightarrow$ "727".

Let's maintain three counters as we iterate through the string from left to right:
\begin{itemize}
    \item \texttt{cnt7}: The number of '7's seen so far.
    \item \texttt{cnt72}: The number of "72" subsequences seen so far.
    \item \texttt{total}: The number of "727" subsequences seen so far (our answer).
\end{itemize}

\textbf{Logic:}
\begin{enumerate}
    \item When we encounter a \textbf{'2'}:
    It can combine with any of the previous \texttt{cnt7} '7's to form a "72".
    So, we update: $\texttt{cnt72} = \texttt{cnt72} + \texttt{cnt7}$.
    
    \item When we encounter a \textbf{'7'}:
    This digit plays two roles.
    \begin{itemize}
        \item First, it can act as the \textit{last} digit ($k$) in a "727" subsequence. It completes any existing "72" subsequences found previously. So, we add the current number of "72"s to our answer: $\texttt{total} = \texttt{total} + \texttt{cnt72}$.
        \item Second, it can act as the \textit{first} digit ($i$) in a future "727" subsequence. So, we increment the count of single '7's: $\texttt{cnt7} = \texttt{cnt7} + 1$.
    \end{itemize}
\end{enumerate}

\textbf{Implementation Note:} When encountering a '7', the order of updates is crucial. We must update \texttt{total} before incrementing \texttt{cnt7}. If we incremented \texttt{cnt7} first, we would imply that a '7' could form a "72" with a future '2' and then complete it itself, which violates the strict index ordering $i < j < k$.

\section*{Complexity Analysis}
\begin{itemize}
    \item \textbf{Time Complexity:} $O(N)$, where $N$ is the length of the string. We iterate through the string exactly once.
    \item \textbf{Space Complexity:} $O(1)$, as we only use three integer variables regardless of input size.
\end{itemize}

\section*{Implementation (C++)}

\begin{lstlisting}[language=C++]
#include <iostream>
#include <string>

using namespace std;

void solve() {
    string s;
    cin >> s;
    
    long long cnt7 = 0;
    long long cnt72 = 0;
    long long total = 0;
    
    for (char c : s) {
        if (c == '7') {
            // 1. This '7' acts as the 3rd digit, completing "72..."
            total += cnt72;
            
            // 2. This '7' acts as the 1st digit for future sequences
            cnt7++;
        } else if (c == '2') {
            // This '2' acts as the 2nd digit, extending "7..."
            cnt72 += cnt7;
        }
    }
    
    cout << total << endl;
}

int main() {
    // Optimization for faster I/O
    ios_base::sync_with_stdio(false);
    cin.tie(NULL);

    int N;
    if (cin >> N) {
        while (N--) {
            solve();
        }
    }
    return 0;
}
\end{lstlisting}

\end{document}
